\documentclass[14pt,a4paper]{extreport}
% preamble.tex — спільна преамбула для всіх лабораторних робіт. Підключається одразу після \documentclass у кожному lab*/main.tex.
%
% ЗБІРКА: LuaLaTeX (не pdfLaTeX). Причина: пакет listings під pdfLaTeX не читає UTF-8 і зупиняє збірку на кириличних коментарях у лістингах програмного коду. LuaLaTeX обробляє UTF-8 нативно.
%
% ШРИФТ: Liberation Serif — метрично сумісний з Times New Roman і містить повний набір кириличних гліфів (TeX Gyre Termes кирилиці не має, через що український текст не відображався взагалі).

\usepackage{fontspec}
\setmainfont{Liberation Serif}
\setmonofont{DejaVu Sans Mono}[Scale=0.82]

\usepackage{polyglossia}
\setmainlanguage{ukrainian}
\setotherlanguage{english}

\usepackage{microtype}

% Поля, інтервал, відступи й заголовки нижче взяті з шаблону НУЛП, за яким написана бакалаврська робота (../bachelor-thesis/thesis/main.tex). Той самий формат приймає В.М. Шиманський і для магістерських, тож усі звіти тримаємо в ньому, а не в тому, що переказували на консультації.
\usepackage[left=25mm,right=15mm,top=15mm,bottom=23mm,footskip=1.5mm]{geometry}
\usepackage{setspace}
% Саме \setstretch{1.5}, а не \onehalfspacing: останній дає близько 1.25.
\setstretch{1.5}

\usepackage{graphicx}
\usepackage{float}
\usepackage{booktabs}
\usepackage{array}
\usepackage{tabularx}
\usepackage{longtable}
\usepackage{enumitem}
% labelindent вирівнює маркер списку з абзацним відступом — так у шаблоні НУЛП.
\setlist{nosep}
\setlist[itemize]{label=\textbullet,leftmargin=*,labelindent=1.27cm}
\setlist[enumerate]{leftmargin=*,labelindent=1.27cm}
\usepackage{indentfirst}
\usepackage{amsmath}

\usepackage{hyperref}
\hypersetup{colorlinks=true,linkcolor=black,urlcolor=blue,citecolor=black}

% xurl дозволяє \url розривати адресу в будь-якому місці, а не лише після скісної риски. Без нього довгий сегмент шляху (напр. «scopus---podrobnaya- instruktsiya-dlya-uchenogo») виходить за межі сторінки, бо дефіс для hyperref не є точкою розриву. Вантажиться після hyperref — перевизначає його \UrlBreaks.
\usepackage{xurl}

% ---------------------------- Лістинги програмного коду ---------------------------- Використано minted (через Pygments), а не listings: listings не вміє коректно обробляти багатобайтові символи й «висмикує» кириличні коментарі з коду, зліплюючи їх в один рядок. Потребує прапорця -shell-escape.
\usepackage{xcolor}
\usepackage[newfloat]{minted}
% labelsep=period дає «Рис. 2.1. Назва» замість «Рис. 2: Назва» — формат НУЛП.
\usepackage[labelsep=period]{caption}

% baselinestretch=1 обов'язковий: без нього лістинг успадковує півторашку основного тексту, і довгий фрагмент коду перестає вміщатися у сторінку. Усередині [H]-флоата це не помилка збірки, а нескінченний цикл порожніх сторінок — lab1 так намолотив 161 тисячу сторінок.
\setminted{
  baselinestretch=1,
  fontsize=\footnotesize,
  linenos=true,
  frame=single,
  framesep=4pt,
  breaklines=true,
  breakanywhere=true,
  tabsize=4,
  numbersep=7pt,
  bgcolor=gray!8,
}
\setminted[python]{python3=true}

\SetupFloatingEnvironment{listing}{name=Лістинг,listname=Перелік лістингів}

% ---------------------------- Нумерація та оформлення заголовків ----------------------------
\usepackage{titlesec}

\setcounter{secnumdepth}{3}
\renewcommand{\thesection}{\arabic{section}}
\renewcommand{\thesubsection}{\arabic{section}.\arabic{subsection}}
\renewcommand{\thesubsubsection}{\arabic{section}.\arabic{subsection}.\arabic{subsubsection}}

% Розділ по центру, 16 пт, великими літерами; слово «РОЗДІЛ» опущене, бо звіт не дисертація. Номер лишається: \numberwithin прив'язує до нього нумерацію рисунків («Рис. 2.1»), і без нумерованих розділів вона б зламалася. \MakeUppercase четвертим аргументом отримує текст заголовка, тому в джерелі заголовки пишуться звичайним регістром.
\titleformat{\section}[block]
  {\centering\bfseries\fontsize{16}{19.2}\selectfont}
  {\thesection.}{0.3em}{\MakeUppercase}
% Підрозділи — зліва з абзацного відступу, звичайним регістром.
\titleformat{\subsection}[block]{\bfseries}{\thesubsection.}{1em}{}
\titleformat{\subsubsection}[block]{\bfseries\itshape}{\thesubsubsection.}{1em}{}

% Незіркові \titlespacing — обов'язкові. Без них titlesec застосовує зірковий варіант, який прибирає абзацний відступ у першого абзацу після заголовка й перекриває indentfirst. Значення — типові для report.
\titlespacing{\section}{0pt}{20pt}{6pt}
\titlespacing{\subsection}{1.27cm}{6pt}{0pt}
\titlespacing{\subsubsection}{1.27cm}{6pt}{0pt}

\setlength{\parindent}{1.27cm}
\setlength{\parskip}{0pt}
\renewcommand{\arraystretch}{1.3}

% Переноси слів вимкнені — так у шаблоні НУЛП. Натомість абзац розганяється міжслівними пробілами, тому tolerance і emergencystretch підняті: без них рядки виходили б за праве поле.
\hyphenpenalty=10000
\exhyphenpenalty=10000
\tolerance=9999
\setlength{\emergencystretch}{3em}

% ---------------------------- Підписи до рисунків і таблиць ----------------------------
\captionsetup[figure]{name=Рис.,justification=centering,font={bf,it}}
\captionsetup[table]{name=Таблиця,justification=raggedleft,singlelinecheck=false,labelfont=bf,textfont=normalfont}

% Рисунки, таблиці й формули нумеруються в межах розділу: «Рис. 2.1», а не «Рис. 5». Так у шаблоні НУЛП, і так само виглядали приклади на консультації.
\numberwithin{equation}{section}
\numberwithin{figure}{section}
\numberwithin{table}{section}

% ---------------------------- Колонтитули ----------------------------
\usepackage{fancyhdr}
\pagestyle{fancy}
\fancyhf{}
\fancyfoot[R]{\thepage}
\renewcommand{\headrulewidth}{0pt}
\fancypagestyle{plain}{%
  \fancyhf{}
  \fancyfoot[R]{\thepage}
  \renewcommand{\headrulewidth}{0pt}
}

% titlepage.tex — титульний аркуш, спільний для всіх дисциплін.
% Потребує визначених у subject.tex команд:
%   \SubjectName, \LecturerName, \StudentName, \StudentGroup, \AcademicYear
%
% Використання у тексті: \MakeLabTitle{1}{Тема роботи}

\newcommand{\MakeLabTitle}[2]{%
\begin{titlepage}
\begin{center}
  {\large\bfseries МІНІСТЕРСТВО ОСВІТИ І НАУКИ УКРАЇНИ}\\[6pt]
  {\large\bfseries Національний університет «Львівська політехніка»}\\[6pt]
  {\large\bfseries Інститут комп'ютерних наук та інформаційних технологій}\\[6pt]
  {\large\bfseries Кафедра «Системи штучного інтелекту»}\\[1.2cm]

  \includegraphics[width=5cm]{../../../../logo}\\[1.2cm]

  {\Large\bfseries ЛАБОРАТОРНА РОБОТА №#1}\\[0.5cm]
  з дисципліни\\[4pt]
  «\SubjectName»\\[0.6cm]
  на тему\\[4pt]
  \textit{#2}\\[2cm]
\end{center}

\begin{flushright}
  {\bfseries Виконав:}\\
  студент групи \StudentGroup\\
  \StudentName\\[1cm]
  {\bfseries Викладач:}\\
  \LecturerName\\
\end{flushright}

\vfill
\begin{center}
  Львів --- \AcademicYear
\end{center}
\end{titlepage}
\clearpage
}


% Subject metadata, read by the shared title page (titlepage.tex).

\newcommand{\SubjectName}{Методологія та процеси моделювання у штучному інтелекті}
% The lecturer's position and degree are not known yet; the other subjects put them on a line above the name, separated by `\\`.
\newcommand{\LecturerName}{Косаревич Р.Я.}
\newcommand{\StudentName}{Галич Максим}
\newcommand{\StudentGroup}{КНСШ-13}
\newcommand{\AcademicYear}{2026}

% Figures are generated artifacts, so they live in .cache/ rather than beside the text (see ARTIFACTS in src/mk1/config.py).
\graphicspath{{../../../../.cache/reports-artifacts/1th term 2026 Autumn/Методологія та процеси моделювання у штучному інтелекті/}}

\begin{document}

\MakeLabTitle{1}{Комплексний дослідницький аналіз даних, просунута попередня обробка та аугментація даних для моделей машинного навчання}

\textbf{Мета роботи.} Навчитися проводити глибинний дослідницький аналіз даних (EDA), виявляти складноструктуровані аномалії, застосовувати методи аугментації та синтезу даних для вирішення проблеми дисбалансу класів чи обмеженого обсягу вибірки.

Усі обчислення виконано мовою Python із бібліотеками scikit-learn, imbalanced-learn, SciPy та pandas. Код організовано як пакет \texttt{mk1}: кожен етап роботи --- окремий модуль, а сценарій \texttt{lab1/run.py} лише запускає етапи по черзі й перевіряє, що кожен записав свої рисунки та метрики. Весь конвеєр відтворюється однією командою з фіксованим зерном генератора випадкових чисел.

\section{Вибір та обґрунтування набору даних}

\subsection{Опис набору даних}

Для роботи обрано набір даних \textbf{Diabetes 130-US Hospitals for Years 1999--2008} з репозиторію UCI Machine Learning Repository (ідентифікатор 296) [\ref{ref:strack}]. Він містить \textbf{101\,766} записів про госпіталізації пацієнтів із цукровим діабетом у 130 лікарнях США за десять років. Кожен запис описано 50 стовпцями: демографія пацієнта, обставини надходження та виписки, тривалість перебування, кількість аналізів, процедур і препаратів, три діагнози в кодах МКХ-9, результати аналізів глюкози та глікованого гемоглобіну (HbA1c), а також 23 стовпці про призначення окремих протидіабетичних препаратів.

Цільова змінна \texttt{readmitted} фіксує, чи повернувся пацієнт до лікарні протягом 30 днів після виписки, пізніше чи не повернувся взагалі. Повторна госпіталізація протягом 30 днів --- поширений показник якості лікування: вона свідчить про передчасну виписку або недостатню підготовку пацієнта до лікування вдома. Тому задачу сформульовано як \textbf{бінарну класифікацію}: клас 1 --- повторна госпіталізація протягом 30 днів, клас 0 --- усі інші випадки.

Набір даних обрано тому, що він природно містить усі труднощі, з якими працює ця лабораторна робота, і жодну з них не довелося створювати штучно (табл.~\ref{tab:requirements}).

\begin{table}[H]
\caption{Відповідність набору даних вимогам завдання}
\label{tab:requirements}
\begin{tabularx}{\linewidth}{|p{5.2cm}|X|}
\hline
\textbf{Вимога або виклик} & \textbf{Як вона проявляється в наборі даних} \\
\hline
Не менше 10 ознак, числові та категоріальні & 46 ознак після очищення: 8 числових лічильників, вік як порядкова ознака та 37 категоріальних \\
\hline
Дисбаланс класів & 8,98~\% спостережень класу 1, співвідношення класів 1~:~10,1 \\
\hline
Пропуски для імп'ютації & пропуски в 9 стовпцях, від 0,01~\% до 48~\%, частина з них прихована за службовими кодами \\
\hline
Високопотужні категоріальні ознаки & коди діагнозів \texttt{diag\_1}--\texttt{diag\_3} (694--756 значень), спеціальність лікаря (70 значень) \\
\hline
Асиметрія та аномалії & лічильники попередніх візитів на 87--93~\% складаються з нулів, асиметрія сягає 21 \\
\hline
Залежні спостереження & 101\,766 госпіталізацій припадають на 71\,518 пацієнтів \\
\hline
\end{tabularx}
\end{table}

Ознаки, що використовуються в роботі, наведено в табл.~\ref{tab:features}. Решту стовпців вилучено під час очищення.

\begin{table}[H]
\caption{Групи ознак та спосіб їх кодування}
\label{tab:features}
\begin{tabularx}{\linewidth}{|p{3.4cm}|X|p{3.6cm}|}
\hline
\textbf{Група} & \textbf{Ознаки} & \textbf{Кодування} \\
\hline
Числові лічильники (8) & \texttt{time\_in\_hospital}, \texttt{num\_lab\_procedures}, \texttt{num\_procedures}, \texttt{num\_medications}, \texttt{number\_outpatient}, \texttt{number\_emergency}, \texttt{number\_inpatient}, \texttt{number\_diagnoses} & масштабування \\
\hline
Порядкова (1) & \texttt{age}: десять вікових інтервалів від [0--10) до [90--100) & порядкове \\
\hline
Низькопотужні категоріальні (15) & \texttt{race}, \texttt{gender}, \texttt{admission\_type\_id}, \texttt{max\_glu\_serum}, \texttt{A1Cresult}, \texttt{change}, \texttt{diabetesMed} та 8 препаратів, призначених щонайменше 1~\% пацієнтів & One-Hot \\
\hline
Високопотужні категоріальні (7) & \texttt{diag\_1}, \texttt{diag\_2}, \texttt{diag\_3}, \texttt{medical\_specialty}, \texttt{discharge\_disposition\_id}, \texttt{admission\_source\_id}, \texttt{payer\_code} & Target або One-Hot (розд.~\ref{sec:final}) \\
\hline
Рідкісні препарати (15) & препарати, призначені менш ніж 1~\% пацієнтів (два --- жодного разу) & лише в лічильниках препаратів (розд.~\ref{sec:features}) \\
\hline
\end{tabularx}
\end{table}

\subsection{Очищення даних}

Перед аналізом виконано чотири кроки очищення. Кожен із них усуває проблему, яка інакше спотворила б і аналіз, і оцінку моделей (табл.~\ref{tab:cleaning}).

\textbf{Приховані пропуски.} Стовпці \texttt{admission\_type\_id}, \texttt{discharge\_disposition\_id} та \texttt{admission\_source\_id} містять числові коди, значення яких описано в довіднику \texttt{IDS\_mapping.csv}. Частина кодів означає «NULL», «Not Available», «Not Mapped» або «Unknown/Invalid», тобто є пропуском, замаскованим під категорію. Такі коди (5, 6, 8; 18, 25, 26; 9, 15, 17, 20, 21 відповідно) замінено на пропуски. Натомість значення «None» у стовпцях \texttt{max\_glu\_serum} та \texttt{A1Cresult} означає «аналіз не проводився» і є повноцінною категорією. Тому під час читання файлу стандартний список пропусків pandas, що містить і рядок «None», вимкнено, а пропуском вважається лише символ «?».

\textbf{Неможлива подія.} Вилучено 2\,423 госпіталізації, що завершилися смертю пацієнта або переведенням у хоспіс. Такого пацієнта не можуть госпіталізувати повторно, тому його мітка визначена самою випискою, а не чимось, чого модель могла б навчитися.

\textbf{Залежні спостереження.} Один пацієнт може мати кілька госпіталізацій з тією самою історією хвороби. Якщо одна його госпіталізація потрапить до навчальної вибірки, а інша --- до тестової, модель фактично бачитиме відповідь. Тому, як і в оригінальному дослідженні [\ref{ref:strack}], для кожного пацієнта залишено лише першу госпіталізацію.

\textbf{Неінформативні стовпці.} Вилучено ідентифікатори \texttt{encounter\_id} і \texttt{patient\_nbr} та вагу пацієнта \texttt{weight}, пропущену в 96,86~\% записів: заповнити таку ознаку означало б її вигадати. Три записи зі значенням статі «Unknown/Invalid» також вилучено.

\begin{table}[H]
\caption{Кількість записів після кожного кроку очищення}
\label{tab:cleaning}
\begin{tabularx}{\linewidth}{|X|r|}
\hline
\textbf{Крок} & \textbf{Записів} \\
\hline
Вихідний набір даних & 101\,766 \\
\hline
Без смертей і переведень у хоспіс & 99\,343 \\
\hline
Лише перша госпіталізація кожного пацієнта & 69\,990 \\
\hline
Без некоректних значень статі & \textbf{69\,987} \\
\hline
\end{tabularx}
\end{table}

\begin{listing}[H]
\inputminted[firstline=173,lastline=211]{python}{../src/mk1/data.py}
\caption{Очищення даних (\texttt{src/mk1/data.py})}
\label{lst:clean}
\end{listing}

Очищений набір поділено на навчальну (80~\%, 55\,989 записів) та тестову (20~\%, 13\,998 записів) вибірки зі стратифікацією за класом. Тестова вибірка використовується лише один раз --- для підсумкової оцінки в розділі~\ref{sec:final}. Дослідницький аналіз і пошук викидів (розділ 2) описують увесь очищений набір. Проте кожне рішення, ухвалене за їхніми результатами, перевірено крос-валідацією лише на навчальній вибірці (розділ~\ref{sec:experiments}), і там само досліджено імп'ютацію, масштабування та аугментацію.

\section{Глибинний дослідницький аналіз}

\subsection{Баланс класів}

Розподіл цільової змінної наведено на рис.~\ref{fig:balance}. У вихідних даних повторну госпіталізацію протягом 30 днів мають 11,2~\% записів. Після очищення частка класу 1 знизилася до \textbf{8,98~\%} (6\,285 з 69\,987), що дає співвідношення класів \textbf{1~:~10,1}. Зниження закономірне: пацієнти з багатьма госпіталізаціями повертаються до лікарні частіше, а від кожного з них тепер залишився один запис.

\begin{figure}[H]
  \centering
  \includegraphics[width=\linewidth]{lab1/class_balance}
  \caption{Розподіл цільової змінної до і після очищення}
  \label{fig:balance}
\end{figure}

Через такий дисбаланс точність (accuracy) непридатна як метрика: модель, яка завжди передбачає клас 0, має точність 91,0~\% і не знаходить жодного пацієнта з ризиком. Тому в роботі використано площу під кривою точність-повнота (PR-AUC) як основну метрику [\ref{ref:saito}], а також ROC-AUC, F1-міру, точність і повноту для класу 1 та оцінку Брієра.

\subsection{Розподіли числових ознак та перевірка нормальності}

Гістограми та Q-Q графіки числових ознак наведено на рис.~\ref{fig:dist} та~\ref{fig:qq}. Найбільш незвичну форму мають лічильники попередніх візитів: \texttt{number\_outpatient}, \texttt{number\_emergency} та \texttt{number\_inpatient} дорівнюють нулю в 87,0~\%, 92,7~\% та 88,3~\% записів відповідно, а решта значень утворює довгий хвіст до 42 візитів. Такий розподіл називають розподілом із надлишком нулів (zero-inflated). Ще 44,1~\% нулів має \texttt{num\_procedures}.

\begin{figure}[H]
  \centering
  \includegraphics[width=\linewidth]{lab1/numeric_distributions}
  \caption{Розподіли числових ознак}
  \label{fig:dist}
\end{figure}

\begin{figure}[H]
  \centering
  \includegraphics[width=\linewidth]{lab1/qq_plots}
  \caption{Q-Q графіки числових ознак відносно стандартного нормального розподілу}
  \label{fig:qq}
\end{figure}

Нормальність кожної ознаки перевірено двома тестами (табл.~\ref{tab:normality}). Тест Шапіро-Уїлка виконано на випадковій підвибірці з 5\,000 записів, бо для більших вибірок SciPy не гарантує точності його p-значення. Тест Д'Агостіно-Пірсона $K^2$ такого обмеження не має і виконаний на всіх 69\,987 записах.

\begin{table}[H]
\caption{Перевірка нормальності числових ознак}
\label{tab:normality}
\begin{tabularx}{\linewidth}{|X|r|r|r|r|r|}
\hline
\textbf{Ознака} & \textbf{Асим.} & \textbf{Ексцес} & \textbf{W} & \textbf{$p_W$} & \textbf{$K^2$} \\
\hline
time\_in\_hospital & 1,18 & 1,01 & 0,883 & $6 \cdot 10^{-52}$ & 11\,903 \\
\hline
num\_lab\_procedures & $-0{,}22$ & $-0{,}30$ & 0,983 & $9 \cdot 10^{-24}$ & 914 \\
\hline
num\_procedures & 1,23 & 0,56 & 0,785 & $6 \cdot 10^{-63}$ & 11\,709 \\
\hline
num\_medications & 1,43 & 3,85 & 0,910 & $2 \cdot 10^{-47}$ & 19\,942 \\
\hline
number\_outpatient & 9,70 & 185 & 0,280 & $3 \cdot 10^{-88}$ & 108\,589 \\
\hline
number\_emergency & 21,23 & 1\,209 & 0,217 & $4 \cdot 10^{-90}$ & 159\,048 \\
\hline
number\_inpatient & 5,59 & 47,3 & 0,328 & $1 \cdot 10^{-86}$ & 76\,256 \\
\hline
number\_diagnoses & $-0{,}72$ & $-0{,}37$ & 0,846 & $1 \cdot 10^{-56}$ & 5\,604 \\
\hline
\end{tabularx}
\end{table}

Обидва тести відхиляють гіпотезу нормальності для всіх восьми ознак: найбільше p-значення тесту $K^2$ становить $3 \cdot 10^{-199}$. Цей результат слід тлумачити обережно. За десятків тисяч спостережень тест відхиляє гіпотезу навіть через незначне відхилення від нормальності, тому про форму розподілу більше говорять Q-Q графіки, асиметрія та ексцес. Навіть майже симетрична \texttt{num\_lab\_procedures} (асиметрія $-0{,}22$) нормальною не є через дискретність значень. Решта ознак відхиляються від нормальності суттєво: Q-Q графіки лічильників візитів мають форму сходинок, а ексцес \texttt{number\_emergency} перевищує 1\,200.

Висновок для подальших кроків: методи, що спираються на нормальність, тут некоректні. Це стосується насамперед Z-оцінки для пошуку викидів. Чи варто наближати розподіли до нормального степеневими перетвореннями, перевірено окремо (розділи~\ref{sec:scaling} та~\ref{sec:experiments}).

\subsection{Кореляційний аналіз і мультиколінеарність}

Матриці кореляцій Пірсона та Спірмена наведено на рис.~\ref{fig:corr}. Обидві матриці майже збігаються: найбільша різниця становить 0,05 для пари \texttt{number\_outpatient}--\texttt{number\_emergency}. Отже, зв'язки між ознаками, де вони є, монотонні та близькі до лінійних. Найсильніші зв'язки: \texttt{num\_medications} з \texttt{time\_in\_hospital} ($r = 0{,}47$) та з \texttt{num\_procedures} ($r = 0{,}40$), \texttt{time\_in\_hospital} з \texttt{num\_lab\_procedures} ($r = 0{,}33$). Зв'язок кожної ознаки з цільовою змінною слабкий: найсильніший має кількість стаціонарних госпіталізацій за попередній рік \texttt{number\_inpatient} ($r = 0{,}10$).

\begin{figure}[H]
  \centering
  \includegraphics[width=\linewidth]{lab1/correlation}
  \caption{Матриці кореляцій Пірсона та Спірмена (числові ознаки, вік і ціль)}
  \label{fig:corr}
\end{figure}

Для кількісної оцінки мультиколінеарності обчислено коефіцієнт збільшення дисперсії $\mathrm{VIF}_j = 1 / (1 - R_j^2)$, де $R_j^2$ --- коефіцієнт детермінації регресії ознаки $j$ на всі інші. Найбільше значення має \texttt{num\_medications}: $\mathrm{VIF} = 1{,}54$, що набагато нижче за поріг 10. Отже, \textbf{вихідні числові ознаки мультиколінеарності не мають}.

Мультиколінеарність у цьому наборі даних зосереджена серед категоріальних ознак. Для них зв'язок виміряно V-коефіцієнтом Крамера з поправкою на зміщення [\ref{ref:bergsma}] (рис.~\ref{fig:cramer}). Найсильніше пов'язані \texttt{change} з \texttt{insulin} ($V = 0{,}61$), \texttt{diabetesMed} з \texttt{insulin} ($V = 0{,}57$) та \texttt{diabetesMed} з \texttt{change} ($V = 0{,}51$). Причина в тому, що \texttt{change} (чи змінювали дозу будь-якого препарату) та \texttt{diabetesMed} (чи призначено будь-який протидіабетичний препарат) є похідними від стовпців окремих препаратів, а інсулін призначено більш ніж половині пацієнтів. З цільовою змінною найсильніше пов'язаний спосіб виписки \texttt{discharge\_disposition\_id} ($V = 0{,}14$): пацієнтів, переведених до інших закладів, госпіталізують повторно частіше.

\begin{figure}[H]
  \centering
  \includegraphics[width=0.9\linewidth]{lab1/cramers_v}
  \caption{Зв'язок між категоріальними ознаками за V-коефіцієнтом Крамера}
  \label{fig:cramer}
\end{figure}

\subsection{Пропущені значення}
\label{sec:missing}

Частку пропусків у кожному стовпці після очищення наведено на рис.~\ref{fig:missing}. Найбільше пропущено спеціальність лікаря \texttt{medical\_specialty} (48,1~\%) та код платника \texttt{payer\_code} (43,5~\%). Приховані пропуски становлять 11,3~\% у \texttt{admission\_type\_id}, 7,2~\% у \texttt{admission\_source\_id} та 4,6~\% у \texttt{discharge\_disposition\_id}. Числові ознаки пропусків не мають.

\begin{figure}[H]
  \centering
  \includegraphics[width=\linewidth]{lab1/missingness}
  \caption{Частка пропусків і частота повторної госпіталізації за наявності та відсутності значення}
  \label{fig:missing}
\end{figure}

Важливо, що сам факт пропуску часто пов'язаний із ціллю (права частина рис.~\ref{fig:missing}, критерій $\chi^2$). Пацієнти без коду платника повертаються до лікарні частіше: 9,9~\% проти 8,3~\% ($\chi^2 = 54{,}6$, $p < 10^{-12}$). Пацієнти без третього діагнозу повертаються майже вдвічі рідше: 5,0~\% проти 9,1~\% ($\chi^2 = 23{,}8$, $p \approx 10^{-6}$), бо менша кількість діагнозів означає легший стан. Пропуск раси також значущий ($p = 0{,}013$), а пропуск спеціальності --- на межі значущості ($p = 0{,}021$). Приховані пропуски, навпаки, з ціллю не пов'язані ($p > 0{,}28$).

Отже, для частини ознак пропуски \textbf{не є випадковими}: факт відсутності значення сам є ознакою. Будь-яке заповнення конкретним значенням цю інформацію знищує. Це визначило вибір методу імп'ютації в розділі~\ref{sec:imputation}.

\subsection{Кардинальність категоріальних ознак}

Кількість унікальних значень кожної категоріальної ознаки наведено на рис.~\ref{fig:card}. Коди діагнозів мають 694--756 значень, а їхній розподіл має довгий хвіст: 10 найчастіших кодів основного діагнозу охоплюють 36,6~\% пацієнтів, 50 найчастіших --- 70,7~\%, а 331 код трапляється менш ніж 10 разів. One-Hot кодування дало б для трьох стовпців діагнозів понад 2\,000 стовпців, більшість яких містять лічені одиниці, і модель не змогла б надійно оцінити їхній вплив. Тому для високопотужних ознак застосовано кодування, що дає один стовпець на ознаку (розділ~\ref{sec:encoding}).

\begin{figure}[H]
  \centering
  \includegraphics[width=\linewidth]{lab1/cardinality}
  \caption{Кардинальність категоріальних ознак і концентрація кодів основного діагнозу}
  \label{fig:card}
\end{figure}

\subsection{Аналіз аномалій та викидів}
\label{sec:outliers}

Для пошуку викидів у числових ознаках порівняно чотири методи:

\begin{itemize}
\item \textbf{міжквартильний інтервал (IQR)}: значення поза межами $[Q_1 - 1{,}5 \cdot \mathrm{IQR};\; Q_3 + 1{,}5 \cdot \mathrm{IQR}]$ хоча б за однією ознакою;
\item \textbf{Z-оцінка}: відхилення від середнього більше ніж на три стандартні відхилення хоча б за однією ознакою;
\item \textbf{Isolation Forest} [\ref{ref:iforest}]: ансамбль із 200 випадкових дерев, де викид ізолюється за меншу кількість розбиттів;
\item \textbf{Local Outlier Factor (LOF)} [\ref{ref:lof}]: порівняння локальної щільності точки зі щільністю її 50 сусідів у стандартизованому просторі.
\end{itemize}

Для Isolation Forest та LOF поріг визначено автоматично (\texttt{contamination="auto"}), щоб кожен метод сам визначив частку викидів. Кількість сусідів LOF збільшено до 50, бо 11~\% записів повністю збігаються з іншим записом за всіма числовими ознаками. За малої кількості сусідів радіус околу такої точки дорівнював би нулю.

\begin{figure}[H]
  \centering
  \includegraphics[width=\linewidth]{lab1/outliers_summary}
  \caption{Частка позначених викидів і частота повторної госпіталізації серед них}
  \label{fig:outliers}
\end{figure}

Методи дали дуже різні результати (рис.~\ref{fig:outliers}): IQR позначив \textbf{33,1~\%} пацієнтів, Z-оцінка --- 9,0~\%, Isolation Forest --- 12,9~\%, LOF --- лише 0,67~\%.

Таку частку для IQR пояснює надлишок нулів. Для \texttt{number\_outpatient} $Q_1 = Q_3 = 0$, тому $\mathrm{IQR} = 0$ і будь-яке ненульове значення стає «викидом». Лише ця ознака дає 13,0~\% позначених записів, а \texttt{number\_inpatient} --- ще 11,7~\% (табл.~\ref{tab:outlier_features}). Z-оцінка позначає менше, але вона спирається на нормальність, яку відхилено в попередньому розділі. До того ж важкий хвіст завищує стандартне відхилення, і частина справжніх екстремумів не перетинає поріг. Для дискретних ознак із надлишком нулів обидва одновимірні методи непридатні.

\begin{table}[H]
\caption{Частка записів, позначених одновимірними методами, за ознаками, \%}
\label{tab:outlier_features}
\begin{tabularx}{\linewidth}{|X|r|r|r|}
\hline
\textbf{Ознака} & \textbf{IQR} & \textbf{Z-оцінка} & \textbf{Нулів} \\
\hline
time\_in\_hospital & 2,00 & 0,92 & 0 \\
\hline
num\_lab\_procedures & 0,14 & 0,04 & 0 \\
\hline
num\_procedures & 5,49 & 0 & 44,07 \\
\hline
num\_medications & 2,67 & 1,49 & 0 \\
\hline
number\_outpatient & 13,03 & 1,81 & 86,97 \\
\hline
number\_emergency & 7,29 & 1,74 & 92,71 \\
\hline
number\_inpatient & 11,71 & 3,42 & 88,29 \\
\hline
number\_diagnoses & 0,34 & 0,34 & 0 \\
\hline
\end{tabularx}
\end{table}

Методи також слабко узгоджуються між собою (рис.~\ref{fig:jaccard}). Найбільший індекс Жаккара мають Z-оцінка та Isolation Forest (0,45). Лише 0,24~\% записів позначили всі чотири методи, а хоча б один --- 33,7~\%. LOF майже не перетинається з іншими методами (індекс не вище 0,03), бо шукає інше: точки з аномально низькою щільністю серед близьких сусідів, а не екстремальні значення. На проєкції на дві головні компоненти (рис.~\ref{fig:outliers_pca}) Z-оцінка та Isolation Forest позначають периферію хмари точок, IQR --- ще й значну частину її ядра, а LOF --- окремі точки, розкидані всередині неї.

\begin{figure}[H]
  \centering
  \includegraphics[width=0.62\linewidth]{lab1/outliers_jaccard}
  \caption{Узгодженість методів пошуку викидів (індекс Жаккара)}
  \label{fig:jaccard}
\end{figure}

\begin{figure}[H]
  \centering
  \includegraphics[width=\linewidth]{lab1/outliers_pca}
  \caption{Позначені викиди на проєкції на дві головні компоненти}
  \label{fig:outliers_pca}
\end{figure}

Головний висновок дає частота повторної госпіталізації серед позначених записів (права частина рис.~\ref{fig:outliers}). Серед викидів за Z-оцінкою вона становить \textbf{14,4~\%} проти 8,4~\% серед решти, за Isolation Forest --- 11,8~\% проти 8,6~\%. Отже, «викиди» --- це не помилки вимірювання, а пацієнти, які часто звертаються по допомогу і саме тому мають підвищений ризик. Їхні значення правдоподібні: 42 візити до приймального відділення за рік трапляються рідко, але можливі. Вилучення таких записів прибрало б із навчальної вибірки найінформативніші приклади класу 1. Тому викиди залишено. Це рішення перевірено експериментально в розділі~\ref{sec:experiments}.

\section{Комплексна обробка та трансформація даних}

\subsection{Заповнення пропусків}
\label{sec:imputation}

\subsubsection{Числові ознаки: експеримент із прихованими значеннями}

Числові ознаки цього набору пропусків не мають, тому методи їх заповнення порівняно на значеннях, які приховано навмисно. Так справжнє значення відоме, і точність кожного методу можна виміряти. На стратифікованій підвибірці з 10\,000 навчальних записів приховано 20~\% значень семи числових ознак за двома механізмами:

\begin{itemize}
\item \textbf{MCAR} (повністю випадкові пропуски): кожне значення приховано з імовірністю 0,2;
\item \textbf{MAR} (випадкові пропуски, що залежать від спостережуваної ознаки): імовірність приховати значення зростає лінійно з рангом тривалості госпіталізації \texttt{time\_in\_hospital}, від 0 для найкоротших до 0,4 для найдовших. Сама \texttt{time\_in\_hospital} лишається повною.
\end{itemize}

Порівняно три методи: заповнення медіаною, \texttt{KNNImputer} (5 сусідів у стандартизованому просторі) та \texttt{IterativeImputer} --- реалізацію MICE [\ref{ref:mice}], що по черзі моделює кожну ознаку регресією на решту (байєсівська гребенева регресія, 15 ітерацій). Якість оцінено трьома показниками на прихованих значеннях:

\begin{itemize}
\item нормована похибка $\mathrm{NRMSE} = \mathrm{RMSE} / \sigma$, де $\sigma$ --- стандартне відхилення ознаки;
\item зміщення: середня різниця між заповненим і справжнім значенням, поділена на $\sigma$;
\item відношення стандартного відхилення заповненої ознаки до справжнього: наскільки метод стискає розкид.
\end{itemize}

\begin{figure}[H]
  \centering
  \includegraphics[width=\linewidth]{lab1/imputation_numeric}
  \caption{Похибка та зміщення методів заповнення числових ознак}
  \label{fig:imp_num}
\end{figure}

\begin{table}[H]
\caption{Середні за сімома ознаками показники методів заповнення}
\label{tab:imp_num}
\begin{tabularx}{\linewidth}{|X|r|r|r|r|r|r|}
\hline
 & \multicolumn{3}{c|}{\textbf{MCAR}} & \multicolumn{3}{c|}{\textbf{MAR}} \\
\cline{2-7}
\textbf{Метод} & \textbf{NRMSE} & \textbf{|зміщ.|} & \textbf{$\sigma$-відн.} & \textbf{NRMSE} & \textbf{|зміщ.|} & \textbf{$\sigma$-відн.} \\
\hline
Медіана & 1,030 & 0,247 & 0,90 & 1,068 & 0,297 & 0,89 \\
\hline
KNNImputer & 1,024 & 0,035 & 0,92 & 1,070 & 0,042 & 0,92 \\
\hline
MICE & \textbf{0,935} & \textbf{0,017} & 0,91 & \textbf{0,974} & \textbf{0,010} & 0,91 \\
\hline
\end{tabularx}
\end{table}

Результати наведено на рис.~\ref{fig:imp_num} і в табл.~\ref{tab:imp_num}. Найменшу похибку має MICE, але навіть у неї NRMSE близька до одиниці. Іншими словами, жоден метод не відновлює приховане значення помітно краще, ніж середнє значення ознаки. Це прямо випливає з кореляційного аналізу: при $\mathrm{VIF} \leq 1{,}54$ ознаки майже не несуть інформації одна про одну, і передбачити одну з них за іншими неможливо.

Методи суттєво різняться \textbf{зміщенням}. Медіана зміщує заповнені значення в середньому на 0,25$\sigma$ навіть за MCAR, бо для асиметричних ознак медіана не збігається із середнім: для лічильників вона менша за нього, для лівоасиметричної \texttt{number\_diagnoses} --- більша. За MAR приховано переважно значення пацієнтів із довгим перебуванням, які мають більше препаратів і процедур. Медіана цього не враховує й занижує, наприклад, \texttt{num\_medications} на 0,45$\sigma$. KNN та MICE використовують \texttt{time\_in\_hospital} і лишаються практично незміщеними (у середньому не більше 0,04$\sigma$). Зміщене заповнення спотворює зв'язок ознаки з ціллю, тому для числових ознак обрано \textbf{MICE}: вона дає найменшу похибку й найменше зміщення, а її вартість лінійна за кількістю рядків. KNN такої ж незміщеності досягає з квадратичною вартістю і без виграшу в точності.

\subsubsection{Категоріальні ознаки}

Категоріальні ознаки мають справжні пропуски. Для них порівняно три стратегії: заповнення модою, KNN-класифікатор (15 сусідів за стандартизованими числовими ознаками) та окрема категорія «Missing». Перші дві перевірено тим самим способом: 20~\% наявних значень приховано й виміряно частку правильно відновлених категорій (рис.~\ref{fig:imp_cat}, табл.~\ref{tab:imp_cat}).

\begin{figure}[H]
  \centering
  \includegraphics[width=\linewidth]{lab1/imputation_categorical}
  \caption{Відновлення прихованих категорій і спотворення розподілу при заповненні модою}
  \label{fig:imp_cat}
\end{figure}

\begin{table}[H]
\caption{Відновлення прихованих категоріальних значень, \%}
\label{tab:imp_cat}
\begin{tabularx}{\linewidth}{|X|r|r|r|r|}
\hline
\textbf{Ознака} & \textbf{Пропусків} & \textbf{Мода} & \textbf{KNN} & \textbf{Частка моди} \\
\hline
race & 2,8 & \textbf{78,4} & 77,0 & 76,9 → 77,6 \\
\hline
payer\_code & 43,5 & \textbf{54,5} & 53,6 & 50,6 → 72,1 \\
\hline
medical\_specialty & 47,4 & 30,7 & \textbf{34,9} & 29,8 → 63,1 \\
\hline
admission\_type\_id & 11,7 & 57,9 & \textbf{63,9} & 57,1 → 62,1 \\
\hline
admission\_source\_id & 7,7 & 59,3 & \textbf{66,6} & 58,0 → 61,2 \\
\hline
discharge\_disposition\_id & 4,6 & \textbf{66,6} & 66,0 & 66,5 → 68,0 \\
\hline
\end{tabularx}
\end{table}

Стовпці «Мода» та «KNN» --- частка правильно відновлених значень, «Частка моди» --- частка найчастішої категорії серед наявних значень і після заповнення всіх пропусків модою. KNN-класифікатор помітно перевершує моду для ознак надходження (на 6--7 відсоткових пунктів) і для спеціальності лікаря (на 4 пункти), але й так правильно відновлює лише 35~\% спеціальностей. Для раси, платника та способу виписки він не кращий за моду. Заповнення модою стовпців із майже половиною пропусків сильно спотворює розподіл: частка внутрішньої медицини серед спеціальностей зростає з 29,8~\% до 63,1~\%, а частка платника Medicare --- з 50,6~\% до 72,1~\%. Половина пацієнтів отримує вигадану спеціальність, і найчастіша категорія штучно стає домінантною.

З урахуванням того, що самі пропуски пов'язані з ціллю (розділ~\ref{sec:missing}), для категоріальних ознак обрано \textbf{окрему категорію «Missing»}. Вона зберігає інформацію про факт пропуску, не вигадує значень і не спотворює розподіл. Усі три стратегії додатково порівняно за якістю моделі в розділі~\ref{sec:experiments}.

\subsection{Кодування категоріальних ознак}
\label{sec:encoding}

Спосіб кодування обрано за типом ознаки (табл.~\ref{tab:features}):

\begin{itemize}
\item \textbf{порядкове кодування} для віку: десять інтервалів мають природний порядок, і One-Hot його б втратив;
\item \textbf{One-Hot} для 15 низькопотужних номінальних ознак (від 2 до 5 значень) і для груп діагнозів, створених у розділі~\ref{sec:features};
\item \textbf{Target Encoding} [\ref{ref:target}] для семи високопотужних ознак: кожна категорія замінюється згладженою частотою класу 1 серед пацієнтів із цією категорією. Рідкісна категорія отримує значення, близьке до загальної частоти класу, тож шум від кодів, що трапляються кілька разів, гаситься автоматично.
\end{itemize}

Target Encoding має ризик витоку: якщо кодувати рядок із врахуванням його власної мітки, модель отримує відповідь у вигляді ознаки. Тому використано перехресне навчання (cross-fitting) \texttt{TargetEncoder} з бібліотеки scikit-learn: навчальні дані діляться на 5 частин, і кожна частина кодується середніми, обчисленими лише на чотирьох інших. Як альтернативи порівняно \textbf{Frequency Encoding} (категорія замінюється своєю часткою в навчальній вибірці, мітка не використовується) та One-Hot, у якому категорії, рідші за 50 записів, об'єднано в одну. Результати порівняння наведено в розділі~\ref{sec:experiments}.

% Listings longer than about 40 lines are set outside a float: an [H] float cannot break across pages, and one taller than a page sends LaTeX into an endless loop of empty pages.
\inputminted[firstline=220,lastline=263]{python}{../src/mk1/preprocessing.py}
\captionof{listing}{Збирання попередньої обробки за типами ознак (\texttt{src/mk1/preprocessing.py})}
\label{lst:prep}

\subsection{Нормалізація та степеневі перетворення}
\label{sec:scaling}

Порівняно чотири перетворення числових ознак: \texttt{StandardScaler}, \texttt{RobustScaler} та степеневі перетворення Yeo-Johnson [\ref{ref:yeojohnson}] і Box-Cox. Box-Cox визначене лише для додатних значень, а лічильники починаються з нуля, тому перед ним кожну ознаку зсунуто так, щоб її мінімум на навчальній вибірці дорівнював 1. Вплив перетворень на форму розподілу наведено на рис.~\ref{fig:scaling} і в табл.~\ref{tab:skew}.

\begin{figure}[H]
  \centering
  \includegraphics[width=\linewidth]{lab1/scaling}
  \caption{Форма розподілу після масштабування та степеневих перетворень}
  \label{fig:scaling}
\end{figure}

\begin{table}[H]
\caption{Асиметрія числових ознак після перетворень}
\label{tab:skew}
\begin{tabularx}{\linewidth}{|X|r|r|r|r|}
\hline
\textbf{Ознака} & \textbf{Вихідна} & \textbf{Yeo-Johnson} & \textbf{Box-Cox} & \textbf{$\lambda_{YJ}$} \\
\hline
time\_in\_hospital & 1,17 & 0,02 & $-0{,}02$ & $-0{,}12$ \\
\hline
num\_lab\_procedures & $-0{,}22$ & $-0{,}23$ & $-0{,}28$ & 0,99 \\
\hline
num\_procedures & 1,24 & 0,18 & 0,18 & $-0{,}52$ \\
\hline
num\_medications & 1,44 & 0,02 & 0,03 & 0,24 \\
\hline
number\_outpatient & 9,62 & 2,20 & 2,20 & $-7{,}35$ \\
\hline
number\_emergency & 17,00 & 3,30 & 3,30 & $-16{,}59$ \\
\hline
number\_inpatient & 5,60 & 2,38 & 2,38 & $-9{,}83$ \\
\hline
number\_diagnoses & $-0{,}73$ & $-0{,}19$ & $-0{,}24$ & 2,15 \\
\hline
\textbf{Середнє |асиметрії|} & 4,63 & \textbf{1,06} & 1,08 & \\
\hline
\end{tabularx}
\end{table}

\texttt{StandardScaler} і \texttt{RobustScaler} є лінійними перетвореннями, тому асиметрію не змінюють взагалі: після них вона збігається з вихідною (стовпець «Вихідна»). \texttt{RobustScaler} для ознак із надлишком нулів до того ж вироджується: їхній $\mathrm{IQR} = 0$, і scikit-learn замінює масштаб на 1, тож перетворення лише віднімає медіану, яка дорівнює нулю.

Степеневі перетворення зменшили середню абсолютну асиметрію з 4,63 до 1,06, а для \texttt{time\_in\_hospital} і \texttt{num\_medications} --- майже до нуля. Для лічильників візитів картина інша. Оцінене $\lambda$ сягає $-16{,}6$, і перетворення відображає всі ненульові значення майже в одну точку (рис.~\ref{fig:scaling}): лічильник фактично стає індикатором «чи був хоча б один візит». Нормальності воно не досягає, бо 88--93~\% маси лишаються в нулі, і тест Д'Агостіно після перетворення так само відхиляє гіпотезу для всіх ознак. Водночас воно знищує інформацію про кількість візитів: пацієнт з одним візитом і пацієнт із двадцятьма стають однаковими.

За критерієм форми розподілу найкращою є \textbf{Yeo-Johnson}. Вона визначена для нулів і від'ємних значень, тому не потребує штучного зсуву, а це важливо для агрегатних ознак із розділу~\ref{sec:features}, центрованих навколо нуля. Тому Yeo-Johnson взято як еталонне перетворення для експериментів. Остаточний вибір зроблено за якістю моделі (розділ~\ref{sec:experiments}), і для лінійної моделі він виявився іншим саме через втрату кількості візитів.

\section{Аугментація та генерування даних}

\subsection{Синтез нових спостережень}
\label{sec:synthesis}

Для подолання дисбалансу порівняно шість методів передискретизації:

\begin{itemize}
\item \textbf{ROS} (random oversampling): випадкове дублювання записів класу 1, базовий варіант;
\item \textbf{SMOTE} [\ref{ref:smote}]: нова точка лежить на відрізку між записом класу 1 та одним із його п'яти найближчих сусідів того ж класу;
\item \textbf{SMOTENC}: варіант SMOTE для змішаних даних. Числові ознаки інтерполюються, а категоріальні отримують найчастіше значення серед сусідів. Працює на вихідних категоріях до кодування;
\item \textbf{ADASYN} [\ref{ref:adasyn}]: як SMOTE, але генерує більше точок біля тих записів класу 1, серед сусідів яких більше записів класу 0, тобто біля межі класів;
\item \textbf{керований гаусів шум}: копії випадкових записів класу 1, до неперервних ознак яких додано шум $\mathcal{N}(0, (\sigma \cdot s_j)^2)$, де $s_j$ --- стандартне відхилення ознаки $j$, а $\sigma$ --- параметр керування. Бінарні One-Hot стовпці копіюються без змін, бо дробова категорія не має сенсу (лістинг~\ref{lst:noise});
\item \textbf{SMOTE-ENN} [\ref{ref:batista}]: SMOTE з подальшим очищенням правилом відредагованих найближчих сусідів (ENN). Видаляється запис будь-якого класу, якщо хоча б один із трьох його найближчих сусідів належить до іншого класу.
\end{itemize}

Як альтернативу без синтезу даних розглянуто \textbf{ваги класів}: помилка на записі класу 1 у функції втрат важить у 10,1 раза більше.

\begin{listing}[H]
\inputminted[firstline=43,lastline=72]{python}{../src/mk1/augmentation.py}
\caption{Аугментація керованим гаусовим шумом (\texttt{src/mk1/augmentation.py})}
\label{lst:noise}
\end{listing}

Усі методи застосовано \textbf{лише до навчальних даних, після розбиття}. Кожен із них є кроком конвеєра imbalanced-learn, який викликає передискретизацію лише під час навчання. Під час крос-валідації синтетичні записи з'являються лише в навчальних блоках, а валідаційний блок і тестова вибірка складаються тільки зі справжніх пацієнтів. Якби аугментацію виконали до розбиття, синтетичні «двійники» пацієнта потрапили б в оцінювальну вибірку, і метрики виявилися б завищеними.

Розмір класів після передискретизації наведено на рис.~\ref{fig:aug_counts}. Усі методи, крім SMOTE-ENN, вирівнюють класи до 50\,961 запису кожен. SMOTE-ENN поводиться інакше: правило ENN видалило \textbf{28\,245 записів класу 0 (55~\%)} і лише 32 записи класу 1, і співвідношення класів змінилося на протилежне, 1~:~2,2 на користь класу 1. Причина в тому, що класи перекриваються настільки сильно, що більшість пацієнтів класу 0 після SMOTE мають серед трьох найближчих сусідів хоча б одну точку класу 1.

\begin{figure}[H]
  \centering
  \includegraphics[width=\linewidth]{lab1/augmentation_counts}
  \caption{Розмір класів навчальної вибірки після передискретизації}
  \label{fig:aug_counts}
\end{figure}

Розташування синтетичних точок показано на рис.~\ref{fig:aug_pca}. На площині двох головних компонент синтетичні точки всіх методів займають ту саму область, що й справжні записи класу 1, а ця область повністю накладається на клас 0. Уже звідси видно, що жоден метод не створює нової розділювальної інформації: він лише переважує ту саму область простору ознак.

\begin{figure}[H]
  \centering
  \includegraphics[width=\linewidth]{lab1/augmentation_pca}
  \caption{Реальні та синтетичні спостереження на проєкції на дві головні компоненти}
  \label{fig:aug_pca}
\end{figure}

Якість синтетичних точок оцінено кількісно двома способами (рис.~\ref{fig:aug_territory}).

\begin{figure}[H]
  \centering
  \includegraphics[width=\linewidth]{lab1/augmentation_territory}
  \caption{Оточення синтетичних точок і можливість відрізнити їх від реальних}
  \label{fig:aug_territory}
\end{figure}

\textbf{Оточення.} Для кожної точки обчислено частку записів класу 0 серед п'яти її найближчих сусідів. Записи класу 1 розділено навпіл: методи генерують точки з першої половини, а сусідів шукають серед класу 0 та другої половини. Без такого розділення найближчим сусідом синтетичної точки здебільшого був би пацієнт, з якого її створено, і всі методи виглядали б кращими, ніж є.

Для випадкової точки ця частка становить 95,3~\%, для справжніх записів класу 1 --- 92,3~\%. Отже, пацієнти з повторною госпіталізацією майже не утворюють окремих скупчень: їхнє оточення лише на 3 відсоткові пункти відрізняється від випадкового. Синтетичні точки опиняються в такому самому оточенні (91,8--93,5~\%). Трохи глибше на територію класу 0 заходять ADASYN (93,3~\%), що за задумом генерує точки біля межі класів, та SMOTENC (93,5~\%). Синтез не створює окремої області класу 1, бо її немає і в реальних даних, а лише відтворює наявне перекриття класів. Саме цю проблему розглядали на лекції на прикладі виявлення шахрайства: інтерполяція між записами меншого класу, перемішаними з більшим, дає точки в зоні перекриття.

\textbf{Відмінність від реальних даних.} Другий спосіб --- змагальна перевірка (adversarial validation). Класифікатор на основі градієнтного бустингу навчено відрізняти 5\,000 синтетичних записів класу 1 від 5\,000 реальних, і його якість виміряно ROC-AUC на крос-валідації. Значення 0,5 означало б, що синтетичні точки неможливо відрізнити від реальних. Насправді SMOTE, ADASYN та гаусів шум розпізнаються \textbf{безпомилково} (ROC-AUC = 1,000), а SMOTENC --- майже безпомилково (0,983).

Причина --- дискретність реальних даних. Усі лічильники цілочислові, а One-Hot стовпці містять лише 0 і 1. SMOTE та ADASYN інтерполюють і бінарні стовпці: 99,3~\% та 99,4~\% їхніх синтетичних записів містять дробове значення на зразок «0,37 інсуліну з підвищеною дозою», якого не має жоден пацієнт. SMOTENC категорії не інтерполює, а лічильники округлює до цілого, але похідні ознаки інтерполює незалежно від лічильників, з яких їх обчислено. У 87~\% його синтетичних записів \texttt{meds\_per\_day} уже не дорівнює відношенню кількості препаратів до тривалості перебування, тоді як у реальних записах ця рівність виконується завжди. Гаусів шум зсуває з цілочислової сітки всі неперервні стовпці. Будь-яка ознака, за якою класифікатор відрізняє синтетичні точки, доступна і моделі, навченій на передискретизованих даних. Модель може навчитися правилу «нецілочислове значення --- клас 1», яке на реальних пацієнтах тестової вибірки ніколи не спрацює. Наслідки цього видно в розділі~\ref{sec:experiments}.

\subsection{Інженерія ознак}
\label{sec:features}

На основі вихідних ознак створено три групи нових ознак (лістинг~\ref{lst:features}):

\begin{itemize}
\item \textbf{доменні ознаки}: загальна кількість візитів за попередній рік \texttt{total\_visits}, інтенсивність лікування \texttt{meds\_per\_day} і \texttt{labs\_per\_day} (препарати й аналізи на день перебування), кількість призначених препаратів \texttt{n\_meds\_used} і препаратів зі зміненою дозою \texttt{n\_meds\_changed} за всіма 23 стовпцями препаратів, а також групи діагнозів \texttt{diag\_*\_group}. Останні зводять 694--756 кодів МКХ-9 до дев'яти груп [\ref{ref:strack}]: хвороби системи кровообігу, органів дихання, органів травлення, кістково-м'язової та сечостатевої систем, діабет, травми, новоутворення та інші;
\item \textbf{агрегації за категоріями}: відхилення кількості препаратів і тривалості перебування від середнього для спеціальності лікаря та для групи основного діагнозу. Ознака показує, наскільки нетиповим є лікування для такого профілю пацієнта. Групові середні обчислюються в методі \texttt{fit} і лише на навчальних даних;
\item \textbf{поліноміальні ознаки}: квадрати та попарні добутки восьми вихідних числових ознак після перетворення Yeo-Johnson, разом 44 ознаки.
\end{itemize}

\begin{listing}[H]
\inputminted[firstline=115,lastline=138]{python}{../src/mk1/features.py}
\caption{Доменні ознаки та агрегації за категоріями (\texttt{src/mk1/features.py})}
\label{lst:features}
\end{listing}

\begin{figure}[H]
  \centering
  \includegraphics[width=\linewidth]{lab1/features}
  \caption{Взаємна інформація ознак із ціллю та VIF після інженерії ознак}
  \label{fig:features}
\end{figure}

Інформативність ознак оцінено взаємною інформацією з цільовою змінною (рис.~\ref{fig:features}, ліворуч). Для цілочислових ознак використано оцінку за таблицею спряженості, бо оцінка за найближчими сусідами не працює на численних однакових значеннях. Чотири з шести найінформативніших ознак --- створені: найвище стоїть відхилення кількості препаратів від середнього для спеціальності (0,0038 нат), далі \texttt{number\_inpatient} (0,0037), \texttt{labs\_per\_day} (0,0035) та відхилення тривалості перебування від середнього для спеціальності (0,0035). Абсолютні значення малі: ентропія цілі становить 0,30 нат, тож кожна ознака окремо пояснює вкрай малу частку невизначеності.

Інженерія ознак створює мультиколінеарність, якої у вихідних даних не було (рис.~\ref{fig:features}, праворуч). \texttt{total\_visits} --- точна сума трьох лічильників, тому VIF усіх чотирьох нескінченний. Агрегатні ознаки мають вигляд $x - \bar{x}_{\text{група}}$, тому майже лінійно залежать від $x$: VIF \texttt{time\_in\_hospital} зріс з 1,39 до 144. Серед 44 поліноміальних ознак медіанний VIF становить 1,45, але шість ознак мають VIF понад 10. Квадрат лічильника візитів після Yeo-Johnson є точною лінійною функцією самого лічильника, бо перетворена ознака має практично лише два значення. Деревоподібним моделям мультиколінеарність не шкодить. Логістична регресія з L2-регуляризацією лишається стійкою, але її коефіцієнти втрачають інтерпретованість.

Групи діагнозів не замінюють, а доповнюють вихідні коди. Зв'язок із ціллю за V Крамера з поправкою на зміщення становить 0,106 для кодів \texttt{diag\_1} і лише 0,030 для їхніх груп: групування втрачає інформацію, а Target Encoding зберігає її в одному стовпці. Тому в конвеєрі використано обидва подання.

\section{Експериментальна перевірка обраних методів}
\label{sec:experiments}

\subsection{Методика}

Висновки розділів 2--4 отримано окремими аналізами: кожен метод оцінено за його власним критерієм. Остаточне обґрунтування вибору дає якість моделі, навченої на оброблених даних. Тому всі методи зібрано в один конвеєр (лістинг~\ref{lst:pipeline}), керований набором параметрів: спосіб імп'ютації, кодування, масштабування, обробка викидів, набір ознак і метод балансування.

\inputminted[firstline=118,lastline=163]{python}{../src/mk1/pipeline.py}
\captionof{listing}{Конвеєр обробки, аугментації та класифікації (\texttt{src/mk1/pipeline.py})}
\label{lst:pipeline}

Порядок кроків конвеєра і є захистом від витоку даних. Кожен крок, що навчається на даних (групові середні, імп'ютер, кодувальники, масштабування, передискретизація), стоїть усередині конвеєра. Під час крос-валідації його заново навчають на навчальних блоках, а до валідаційного блоку лише застосовують. Передискретизація з бібліотеки imbalanced-learn виконується тільки під час навчання, тож оцінювання завжди відбувається на реальних пацієнтах.

\textbf{Еталонна конфігурація} відповідає висновкам попередніх розділів: категорія «Missing» для пропусків, Target Encoding, перетворення Yeo-Johnson, викиди залишено, доменні та агрегатні ознаки, без балансування. Кожен експеримент змінює лише один її параметр, тому різниця в метриках є наслідком саме цього параметра.

Кожну конфігурацію оцінено двома моделями з різною чутливістю до обробки даних:

\begin{itemize}
\item \textbf{логістична регресія} (L2-регуляризація, $C = 1$): лінійна модель, на яку прямо впливають кодування, масштаб і форма розподілу ознак;
\item \textbf{гістограмний градієнтний бустинг} (250 дерев, крок навчання 0,05, до 31 листка в дереві): деревоподібна модель, інваріантна до монотонних перетворень ознак.
\end{itemize}

Якість оцінено стратифікованою 5-блоковою крос-валідацією на навчальній вибірці (55\,989 записів) і подано як середнє $\pm$ стандартне відхилення між блоками. Основна метрика --- PR-AUC; базовий рівень випадкового класифікатора для неї дорівнює частці класу 1, тобто 0,090. Додатково наведено ROC-AUC, F1-міру при стандартному порозі 0,5, найкращу досяжну F1-міру (за оптимального для блоку порога) та оцінку Брієра, що вимірює якість самих імовірностей.

\subsection{Обробка даних та інженерія ознак}

Результати всіх порівнянь наведено на рис.~\ref{fig:experiments}, а для методів обробки даних та інженерії ознак --- також у табл.~\ref{tab:cv_prep}. Жирним виділено найкращий варіант у межах експерименту для кожної моделі.

\begin{figure}[H]
  \centering
  \includegraphics[width=\linewidth]{lab1/experiments}
  \caption{PR-AUC на 5-блоковій крос-валідації для кожного експерименту}
  \label{fig:experiments}
\end{figure}

\begin{table}[H]
\caption{Крос-валідація методів обробки даних та інженерії ознак}
\label{tab:cv_prep}
\small
\renewcommand{\arraystretch}{1.15}
\begin{tabularx}{\linewidth}{|X|c|c|c|c|}
\hline
 & \multicolumn{2}{c|}{\textbf{Логістична регресія}} & \multicolumn{2}{c|}{\textbf{Градієнтний бустинг}} \\
\cline{2-5}
\textbf{Варіант} & \textbf{PR-AUC} & \textbf{ROC-AUC} & \textbf{PR-AUC} & \textbf{ROC-AUC} \\
\hline
\multicolumn{5}{|l|}{\textit{Імп'ютація категоріальних ознак}} \\
\hline
категорія Missing & \textbf{0,167 ± 0,003} & 0,650 & \textbf{0,179 ± 0,005} & 0,655 \\
\hline
мода & 0,166 ± 0,003 & 0,649 & 0,176 ± 0,005 & 0,654 \\
\hline
KNN & 0,167 ± 0,003 & 0,649 & 0,176 ± 0,004 & 0,653 \\
\hline
\multicolumn{5}{|l|}{\textit{Кодування високопотужних ознак}} \\
\hline
Target & \textbf{0,167 ± 0,003} & 0,650 & 0,179 ± 0,005 & 0,655 \\
\hline
Frequency & 0,145 ± 0,005 & 0,633 & 0,180 ± 0,005 & 0,647 \\
\hline
One-Hot & 0,162 ± 0,002 & 0,642 & \textbf{0,186 ± 0,005} & 0,656 \\
\hline
\multicolumn{5}{|l|}{\textit{Перетворення числових ознак}} \\
\hline
Yeo-Johnson & 0,167 ± 0,003 & 0,650 & 0,179 ± 0,005 & 0,655 \\
\hline
Box-Cox & 0,166 ± 0,003 & 0,650 & 0,179 ± 0,005 & 0,655 \\
\hline
Robust & \textbf{0,172 ± 0,005} & 0,651 & 0,179 ± 0,005 & 0,655 \\
\hline
Standard & \textbf{0,172 ± 0,005} & 0,651 & 0,179 ± 0,005 & 0,655 \\
\hline
\multicolumn{5}{|l|}{\textit{Обробка викидів}} \\
\hline
залишити & \textbf{0,167 ± 0,003} & 0,650 & \textbf{0,179 ± 0,005} & 0,655 \\
\hline
вилучити (Isolation Forest) & 0,163 ± 0,004 & 0,645 & 0,169 ± 0,007 & 0,647 \\
\hline
вилучити (IQR) & 0,149 ± 0,003 & 0,627 & 0,150 ± 0,004 & 0,626 \\
\hline
\multicolumn{5}{|l|}{\textit{Інженерія ознак}} \\
\hline
лише вихідні & 0,165 ± 0,003 & 0,651 & 0,180 ± 0,005 & 0,656 \\
\hline
+ доменні & 0,166 ± 0,003 & 0,650 & \textbf{0,181 ± 0,005} & 0,657 \\
\hline
+ агрегати & \textbf{0,167 ± 0,003} & 0,650 & 0,179 ± 0,005 & 0,655 \\
\hline
+ поліноміальні & 0,167 ± 0,003 & 0,649 & 0,177 ± 0,005 & 0,653 \\
\hline
\end{tabularx}
\end{table}

Абсолютний рівень якості помірний: PR-AUC 0,17--0,19 проти 0,090 у випадкового класифікатора, ROC-AUC близько 0,65. Це узгоджується з результатами аналізу: окремі ознаки слабко пов'язані з ціллю, а пацієнти класу 1 майже не відокремлені від класу 0 у просторі ознак. Тому відмінності між методами невеликі, і їх слід порівнювати зі стандартним відхиленням між блоками (0,002--0,007).

\textbf{Імп'ютація.} Окрема категорія «Missing» найкраща для обох моделей. Для бустингу перевага над модою та KNN становить 0,003, тобто приблизно 0,6 стандартного відхилення, і узгоджується з розділом~\ref{sec:missing}: дерево може безпосередньо використати категорію «Missing», яка несе інформацію про ціль. KNN нічого не додає і в 2,5--3 рази сповільнює навчання конвеєра.

\textbf{Кодування.} Тут моделі розходяться. Для логістичної регресії найкраще Target Encoding (0,167), гірше One-Hot (0,162) і найгірше Frequency Encoding (0,145). Частота категорії не пов'язана з ризиком монотонно, тож лінійна модель не може нею скористатися. Target Encoding дає саме ризик, тобто монотонну шкалу, яку лінійна модель використовує напряму. Для бустингу найкраще One-Hot (0,186): дерево розбиває дані за окремими кодами й не потребує готової шкали. Ціна --- 553 стовпці замість 113 і приблизно втричі довше навчання.

\textbf{Масштабування.} Для бустингу всі чотири перетворення дали однаковий результат (0,179 з точністю до четвертого знака), що експериментально підтверджує інваріантність дерев до монотонних перетворень. Для логістичної регресії лінійне масштабування (Standard, Robust: 0,172) \textbf{перевершило} степеневі перетворення (Yeo-Johnson, Box-Cox: 0,167). Це наслідок ефекту з розділу~\ref{sec:scaling}: перетворення, оптимальне за критерієм нормальності, перетворило лічильники візитів на індикатори і знищило інформацію про кількість попередніх госпіталізацій, найсильнішу числову ознаку. Наближення розподілу до нормального не є метою саме по собі. Перевагу віддано \texttt{RobustScaler}: якість та сама, що в \texttt{StandardScaler}, а для ознак з ненульовим IQR він менш чутливий до екстремальних значень.

\textbf{Викиди.} Вилучення викидів погіршило обидві моделі: Isolation Forest --- на 0,004 і 0,011, IQR --- на 0,018 і 0,029. Це прямо підтверджує висновок розділу~\ref{sec:outliers}: «викиди» тут --- пацієнти з високим ризиком, і вилучення третини вибірки за правилом IQR позбавляє модель найінформативніших прикладів.

\textbf{Інженерія ознак.} Ефект створених ознак не перевищує одного стандартного відхилення. Доменні ознаки трохи покращують обидві моделі. Агрегації за категоріями корисні для логістичної регресії, але не для бустингу, бо дерева будують такі взаємодії самі. Поліноміальні ознаки нічого не дали лінійній моделі й трохи погіршили бустинг, додавши 36 переважно шумових стовпців. Висока взаємна інформація окремих створених ознак (розділ~\ref{sec:features}) не означає нової інформації для моделі: вона здебільшого вже міститься у вихідних ознаках, про що свідчить і високий VIF.

\subsection{Балансування класів}

Результати методів балансування за якістю ранжування наведено в табл.~\ref{tab:cv_resampling}, а за якістю рішень та ймовірностей --- у табл.~\ref{tab:cv_decisions}.

\begin{table}[H]
\caption{Крос-валідація методів балансування класів: якість ранжування}
\label{tab:cv_resampling}
\small
\renewcommand{\arraystretch}{1.15}
\begin{tabularx}{\linewidth}{|X|c|c|c|c|}
\hline
 & \multicolumn{2}{c|}{\textbf{Логістична регресія}} & \multicolumn{2}{c|}{\textbf{Градієнтний бустинг}} \\
\cline{2-5}
\textbf{Метод} & \textbf{PR-AUC} & \textbf{ROC-AUC} & \textbf{PR-AUC} & \textbf{ROC-AUC} \\
\hline
без балансування & \textbf{0,167 ± 0,003} & 0,650 & 0,179 ± 0,005 & 0,655 \\
\hline
ваги класів & 0,166 ± 0,003 & 0,651 & 0,174 ± 0,004 & 0,650 \\
\hline
ROS & 0,165 ± 0,004 & 0,651 & 0,173 ± 0,007 & 0,650 \\
\hline
SMOTE & 0,158 ± 0,005 & 0,639 & 0,126 ± 0,022 & 0,573 \\
\hline
SMOTENC & 0,107 ± 0,002 & 0,550 & 0,121 ± 0,002 & 0,592 \\
\hline
SMOTENC + One-Hot & 0,110 ± 0,004 & 0,552 & 0,124 ± 0,004 & 0,597 \\
\hline
ADASYN & 0,155 ± 0,005 & 0,637 & 0,117 ± 0,015 & 0,551 \\
\hline
гаусів шум ($\sigma = 0{,}1$) & 0,166 ± 0,003 & 0,651 & \textbf{0,180 ± 0,005} & 0,657 \\
\hline
SMOTE-ENN & 0,159 ± 0,004 & 0,647 & 0,153 ± 0,007 & 0,630 \\
\hline
\end{tabularx}
\end{table}

\begin{table}[H]
\caption{Крос-валідація методів балансування класів: рішення та ймовірності}
\label{tab:cv_decisions}
\small
\renewcommand{\arraystretch}{1.15}
\begin{tabularx}{\linewidth}{|X|c|c|c|c|c|c|}
\hline
 & \multicolumn{3}{c|}{\textbf{Логістична регресія}} & \multicolumn{3}{c|}{\textbf{Градієнтний бустинг}} \\
\cline{2-7}
\textbf{Метод} & \textbf{F1\textsubscript{0,5}} & \textbf{F1\textsubscript{max}} & \textbf{Брієр} & \textbf{F1\textsubscript{0,5}} & \textbf{F1\textsubscript{max}} & \textbf{Брієр} \\
\hline
без балансування & 0,006 & 0,233 & 0,079 & 0,012 & 0,236 & 0,079 \\
\hline
ваги класів & 0,219 & 0,233 & 0,229 & 0,228 & 0,236 & 0,193 \\
\hline
ROS & 0,219 & 0,233 & 0,229 & 0,225 & 0,232 & 0,190 \\
\hline
SMOTE & 0,213 & 0,225 & 0,228 & 0,137 & 0,188 & 0,142 \\
\hline
SMOTENC & 0,148 & 0,174 & 0,149 & 0,097 & 0,195 & 0,101 \\
\hline
SMOTENC + One-Hot & 0,148 & 0,175 & 0,143 & 0,085 & 0,197 & 0,097 \\
\hline
ADASYN & 0,213 & 0,222 & 0,226 & 0,148 & 0,175 & 0,161 \\
\hline
гаусів шум ($\sigma = 0{,}1$) & 0,220 & 0,233 & 0,229 & 0,014 & 0,238 & 0,079 \\
\hline
SMOTE-ENN & 0,186 & 0,229 & 0,418 & 0,206 & 0,216 & 0,234 \\
\hline
\end{tabularx}
\end{table}

\textbf{Жоден метод балансування не покращив якість ранжування.} Ваги класів і ROS знизили PR-AUC незначно, синтетичні методи --- суттєво. Для бустингу SMOTE та ADASYN призвели до обвалу: PR-AUC 0,126 та 0,117, ROC-AUC 0,57 та 0,55, з великим розкидом між блоками. Це прямий наслідок результату змагальної перевірки з розділу~\ref{sec:synthesis}. Синтетичні точки SMOTE та ADASYN містять дробові значення бінарних стовпців, а дерево легко будує правило «дробове значення --- клас 1». На навчальних даних це правило бездоганне, на реальних пацієнтах не спрацьовує ніколи. Лінійна модель такого правила побудувати не може, тому постраждала менше (0,158 та 0,155).

Найгірший результат має SMOTENC: PR-AUC 0,107 для логістичної регресії, лише на 0,017 вище за випадковий класифікатор. Спочатку здавалося, що причина у взаємодії з Target Encoding: синтетичні копії рідкісних кодів діагнозів могли б підробити їхнє кодування. Варіант SMOTENC + One-Hot цю гіпотезу спростував, бо дав такий самий результат (0,110). Справжню причину видно з порівняння синтетичних записів із реальними записами класу 1 (табл.~\ref{tab:smotenc}). SMOTENC присвоює синтетичному запису найчастішу категорію серед його сусідів, і найчастіші категорії стають ще частішими. Крім того, imbalanced-learn зберігає цілочисловий тип лічильників і відкидає дробову частину інтерпольованих значень. Через це середня кількість попередніх госпіталізацій синтетичного «пацієнта з повторною госпіталізацією» вдвічі менша, ніж у реальних, і дорівнює рівню класу 0. Модель отримує десятки тисяч записів класу 1, що виглядають як типові пацієнти класу 0, і майже не бачить справжніх.

\begin{table}[H]
\caption{Реальні та синтетичні (SMOTENC) записи класу 1}
\label{tab:smotenc}
\begin{tabularx}{\linewidth}{|X|l|r|r|}
\hline
\textbf{Ознака} & \textbf{Значення} & \textbf{Реальні} & \textbf{SMOTENC} \\
\hline
discharge\_disposition\_id & виписка додому & 48,6~\% & 64,4~\% \\
\hline
medical\_specialty & Missing & 49,8~\% & 75,7~\% \\
\hline
payer\_code & Missing & 47,7~\% & 62,1~\% \\
\hline
diag\_1\_group & хвороби кровообігу & 33,3~\% & 49,0~\% \\
\hline
A1Cresult & аналіз не проводився & 83,0~\% & 97,3~\% \\
\hline
race & Caucasian & 76,6~\% & 94,1~\% \\
\hline
number\_inpatient & середнє значення & 0,363 & 0,166 \\
\hline
\end{tabularx}
\end{table}

Для порівняння: середнє \texttt{number\_inpatient} у класі 0 становить 0,158.

\textbf{Рішення та ймовірності.} Без балансування F1-міра при порозі 0,5 майже нульова (0,006 та 0,012). Модель, навчена на вибірці з 9~\% класу 1, рідко видає ймовірність понад 0,5, тож майже нікого не відносить до класу 1. Ваги класів і ROS піднімають F1\textsubscript{0,5} до 0,22, але найкраща досяжна F1-міра лишається тією самою: 0,233 для логістичної регресії та 0,236 для бустингу. Отже, \textbf{єдиний ефект цих методів --- зсув порога}, і того самого можна досягти без балансування, просто обравши поріг [\ref{ref:elor}]. Зсув має ціну: оцінка Брієра зростає з 0,079 до 0,19--0,23, а для SMOTE-ENN, що змінив співвідношення класів на протилежне, --- до 0,418. Імовірності моделі перестають бути ймовірностями: модель, навчена на вибірці 1~:~1, вважає повторну госпіталізацію подією з шансом близько 50~\%, а не 9~\%. Для клінічної моделі ризику відкалібрована ймовірність є самостійним результатом, і балансування її руйнує.

\textbf{Керований гаусів шум.} Вплив масштабу шуму $\sigma$ наведено на рис.~\ref{fig:noise}. Від $\sigma = 0{,}05$ до $\sigma = 1$ PR-AUC майже не змінюється: від 0,166 до 0,163 для логістичної регресії та від 0,179 до 0,178 для бустингу. Для логістичної регресії шум діє так само, як ROS: зсуває поріг (F1\textsubscript{0,5} = 0,22) і погіршує калібрування. Для бустингу F1\textsubscript{0,5} лишається близькою до нуля за будь-якого $\sigma$, а оцінка Брієра --- такою самою, як без балансування (0,079). Дерева відокремили синтетичні копії, які розпізнаються зі 100-відсотковою точністю, в окремі листки, і для реальних пацієнтів модель лишилася такою самою, як без аугментації.

\begin{figure}[H]
  \centering
  \includegraphics[width=\linewidth]{lab1/noise_sweep}
  \caption{Вплив масштабу гаусового шуму на якість ранжування та рішень}
  \label{fig:noise}
\end{figure}

Висновок для цього набору даних: \textbf{синтетичну аугментацію не застосовано}. Дисбаланс класів враховано вибором порога рішення на прогнозах крос-валідації навчальної вибірки (розділ~\ref{sec:final}). Це дає таку саму F1-міру, як найкращі методи балансування, зберігає найвищу PR-AUC і калібровані ймовірності та не потребує синтетичних даних.

\section{Підсумковий конвеєр і перевірка на тестовій вибірці}
\label{sec:final}

\subsection{Обрані методи та їх обґрунтування}

Методи, обрані за результатами аналізу та експериментів, зведено в табл.~\ref{tab:choices}. Оскільки моделі по-різному реагують на кодування та масштабування, для кожної з них зібрано власну підсумкову конфігурацію.

% A longtable rather than a table float: eight rows of prose run longer than a page, and an [H] float taller than a page sends LaTeX into an endless loop of empty pages.
{\small
\renewcommand{\arraystretch}{1.15}
\begin{longtable}{|p{3.1cm}|p{3.9cm}|p{8.6cm}|}
\caption{Обрані методи обробки та аугментації даних}
\label{tab:choices} \\
\hline
\textbf{Етап} & \textbf{Обраний метод} & \textbf{Обґрунтування} \\
\hline
\endfirsthead
\hline
\textbf{Етап} & \textbf{Обраний метод} & \textbf{Обґрунтування} \\
\hline
\endhead
Пропуски категоріальних ознак & окрема категорія «Missing» & факт пропуску пов'язаний із ціллю ($\chi^2 = 54{,}6$ для коду платника); мода подвоює частку найчастішої категорії; найвища PR-AUC для обох моделей \\
\hline
Пропуски числових ознак & MICE (\texttt{IterativeImputer}) & найменші похибка (NRMSE 0,935) та зміщення (0,017$\sigma$) в експерименті з прихованими значеннями; медіана за MAR зміщує значення до 0,45$\sigma$ \\
\hline
Високопотужні категоріальні ознаки & Target Encoding з перехресним навчанням (лог. регресія); One-Hot з об'єднанням рідкісних категорій (бустинг) & 694--756 кодів діагнозів; лінійній моделі потрібна монотонна шкала ризику (0,167 проти 0,145 у Frequency); дерева краще використовують окремі коди (0,186 проти 0,179) \\
\hline
Інші категоріальні ознаки & One-Hot; вік --- порядкове & 2--5 значень; вікові інтервали мають природний порядок \\
\hline
Числові ознаки & \texttt{RobustScaler} (лог. регресія); для бустингу не впливає & степеневі перетворення перетворюють лічильники візитів на індикатори (0,172 проти 0,167); дерева інваріантні до монотонних перетворень \\
\hline
Викиди & залишено & це пацієнти з підвищеним ризиком (14,4~\% проти 8,4~\%); вилучення погіршує PR-AUC на 0,004--0,029 \\
\hline
Інженерія ознак & доменні ознаки й агрегації (лог. регресія); доменні ознаки (бустинг) & невеликий, але стабільний приріст для кожної моделі; поліноміальні ознаки не допомогли \\
\hline
Дисбаланс класів & без синтезу даних; поріг рішення, обраний на крос-валідації & жоден метод не покращив PR-AUC; синтетичні точки розпізнаються з ROC-AUC 0,98--1,00 і шкодять; ваги класів та ROS лише зсувають поріг і руйнують калібрування \\
\hline
\end{longtable}
}

\subsection{Перехресна перевірка комбінацій}

Експерименти змінювали по одному параметру, тому взаємодію між обраними методами вони не враховують. Тож підсумкові конфігурації ще раз перевірено крос-валідацією до використання тестової вибірки (табл.~\ref{tab:cv_final}). Для порівняння наведено «наївну» конфігурацію, яку часто застосовують без аналізу: заповнення модою, One-Hot для всіх категоріальних ознак, \texttt{StandardScaler}, без створених ознак і без балансування.

\begin{table}[H]
\caption{Крос-валідація підсумкових конфігурацій}
\label{tab:cv_final}
\small
\renewcommand{\arraystretch}{1.15}
\begin{tabularx}{\linewidth}{|X|c|c|c|c|}
\hline
 & \multicolumn{2}{c|}{\textbf{Логістична регресія}} & \multicolumn{2}{c|}{\textbf{Градієнтний бустинг}} \\
\cline{2-5}
\textbf{Конфігурація} & \textbf{PR-AUC} & \textbf{ROC-AUC} & \textbf{PR-AUC} & \textbf{ROC-AUC} \\
\hline
еталонна & 0,167 ± 0,003 & 0,650 & 0,179 ± 0,005 & 0,655 \\
\hline
підсумкова для логістичної регресії & \textbf{0,172 ± 0,005} & 0,651 & 0,179 ± 0,005 & 0,655 \\
\hline
підсумкова для бустингу & 0,161 ± 0,003 & 0,641 & \textbf{0,186 ± 0,008} & 0,656 \\
\hline
наївна & 0,166 ± 0,003 & 0,641 & 0,180 ± 0,006 & 0,652 \\
\hline
\end{tabularx}
\end{table}

Покращення окремих параметрів поєдналися: PR-AUC логістичної регресії зросла з 0,167 до 0,172, бустингу --- з 0,179 до 0,186. Кожна підсумкова конфігурація добра лише для своєї моделі: логістична регресія на конфігурації бустингу дає 0,161, гірше за наївну. Отже, обробку даних слід обирати разом із моделлю, а не окремо від неї.

\subsection{Результати на тестовій вибірці}

Поріг рішення для кожного конвеєра обрано так, щоб максимізувати F1-міру на прогнозах 5-блокової крос-валідації навчальної вибірки. Потім конвеєр навчено на всій навчальній вибірці й один раз застосовано до 13\,998 пацієнтів тестової вибірки, з яких 1\,257 були повторно госпіталізовані (табл.~\ref{tab:test}, рис.~\ref{fig:final}).

\begin{table}[H]
\caption{Результати на тестовій вибірці}
\label{tab:test}
\small
\renewcommand{\arraystretch}{1.15}
\begin{tabularx}{\linewidth}{|X|c|c|c|c|c|c|c|}
\hline
\textbf{Конвеєр} & \textbf{PR-AUC} & \textbf{ROC-AUC} & \textbf{Поріг} & \textbf{F1} & \textbf{Точн.} & \textbf{Повн.} & \textbf{Брієр} \\
\hline
підсумковий, лог. регресія & 0,178 & 0,656 & 0,103 & 0,231 & 0,157 & 0,434 & 0,0789 \\
\hline
наївний, лог. регресія & 0,169 & 0,650 & 0,120 & 0,235 & 0,172 & 0,373 & 0,0794 \\
\hline
підсумковий, бустинг & \textbf{0,187} & 0,652 & 0,134 & \textbf{0,244} & 0,205 & 0,302 & 0,0786 \\
\hline
наївний, бустинг & 0,186 & 0,653 & 0,105 & 0,229 & 0,154 & 0,450 & 0,0787 \\
\hline
\end{tabularx}
\end{table}

\begin{figure}[H]
  \centering
  \includegraphics[width=\linewidth]{lab1/final_curves}
  \caption{Криві точність-повнота та ROC-криві на тестовій вибірці}
  \label{fig:final}
\end{figure}

Тестова вибірка підтвердила результати крос-валідації для логістичної регресії: PR-AUC підсумкового конвеєра 0,178 проти 0,169 у наївного. Для бустингу підсумковий і наївний конвеєри майже рівні (0,187 проти 0,186): перевага 0,006, отримана на крос-валідації, не перевищувала розкиду між блоками і на тестовій вибірці не відтворилася.

Наївний конвеєр виявився сильним суперником, бо в головному він збігається з підсумковим: не вилучає викидів, не синтезує даних, масштабує ознаки лінійно, а для бустингу ще й використовує One-Hot. Найбільший внесок ретельної обробки в цьому наборі даних --- \textbf{уникнення шкідливих кроків}. Вилучення викидів за IQR коштувало до 0,029 PR-AUC, SMOTE і ADASYN для бустингу --- до 0,062, SMOTENC --- до 0,060, Frequency Encoding для лінійної моделі --- 0,022. Ці втрати в кілька разів перевищують виграш від найкращих рішень (0,005--0,007).

Підсумковий конвеєр із бустингом за обраного порога 0,134 відносить до групи ризику 1\,845 пацієнтів (13,2~\% тестової вибірки). Серед них 379 справді повернулися до лікарні протягом 30 днів, тобто точність становить 20,5~\%, у 2,3 раза більше за базову частоту 9,0~\%. Модель знаходить 30,2~\% усіх повторних госпіталізацій. Логістична регресія з порогом 0,103 відбирає ширшу групу (24,8~\% пацієнтів) і знаходить 43,4~\% повторних госпіталізацій за точності 15,7~\%. F1-міра на тестовій вибірці (0,244 та 0,231) близька до найкращої досяжної на крос-валідації для цих конфігурацій (0,238 та 0,233), отже поріг, обраний на навчальних даних, добре переноситься на нові дані.

Імовірності моделей лишилися каліброваними: оцінка Брієра 0,079 нижча, ніж у сталого прогнозу базової частоти ($0{,}090 \cdot 0{,}910 = 0{,}082$). Після будь-якого методу балансування вона становила б 0,19--0,23 (табл.~\ref{tab:cv_decisions}).

\section{Висновки}

У лабораторній роботі виконано комплексний дослідницький аналіз, попередню обробку та аугментацію набору даних Diabetes 130-US Hospitals для задачі передбачення повторної госпіталізації протягом 30 днів. Кожен вибір методу перевірено кількісно.

\begin{enumerate}
\item Обраний набір даних відповідає всім вимогам завдання: після очищення 69\,987 незалежних пацієнтів, 46 числових і категоріальних ознак, дисбаланс класів 1~:~10,1, пропуски в 9 стовпцях і категоріальні ознаки до 756 значень. Очищення усунуло три приховані проблеми: пропуски, замасковані під службові коди; госпіталізації, після яких повторна неможлива; залежні госпіталізації одного пацієнта.

\item Жодна числова ознака не має нормального розподілу: тести Шапіро-Уїлка та Д'Агостіно відхилили гіпотезу з $p < 10^{-23}$, а лічильники візитів на 87--93~\% складаються з нулів. Кореляції Пірсона та Спірмена майже збігаються (максимум $r = 0{,}47$), VIF не перевищує 1,54, тож числові ознаки мультиколінеарності не мають. Вона зосереджена серед категоріальних ознак, похідних від стовпців препаратів (V Крамера до 0,61). Пропуски в коді платника та діагнозах пов'язані з ціллю, тобто не є випадковими.

\item Чотири методи пошуку викидів позначили від 0,7~\% до 33,1~\% записів і слабко узгоджуються між собою (індекс Жаккара не вище 0,45). IQR і Z-оцінка непридатні для дискретних ознак із надлишком нулів. Позначені записи --- це пацієнти з підвищеним ризиком (14,4~\% повторних госпіталізацій проти 8,4~\%), тому їх залишено; вилучення погіршувало якість моделей на 0,004--0,029 PR-AUC.

\item Для категоріальних пропусків обрано окрему категорію «Missing»: вона зберігає інформацію про сам факт пропуску й дала найкращу якість обох моделей. Для числових ознак у разі пропусків найкращим є MICE: NRMSE 0,935 і зміщення до 0,02$\sigma$ проти 0,25--0,30$\sigma$ у медіани. Можливості будь-якого імп'ютера обмежені слабкими зв'язками між ознаками.

\item Високопотужні ознаки закодовано Target Encoding з перехресним навчанням для логістичної регресії та One-Hot з об'єднанням рідкісних категорій для бустингу. Frequency Encoding непридатне для лінійної моделі (PR-AUC 0,145 проти 0,167).

\item Степеневі перетворення Yeo-Johnson і Box-Cox зменшили середню асиметрію з 4,63 до 1,06, але перетворили лічильники візитів на індикатори й погіршили лінійну модель. Обрано \texttt{RobustScaler}; для бустингу всі перетворення дали однаковий результат.

\item Жоден із шести методів балансування (ROS, SMOTE, SMOTENC, ADASYN, керований гаусів шум, SMOTE-ENN) не покращив якість ранжування. Синтетичні записи розпізнаються класифікатором з ROC-AUC 0,98--1,00, бо реальні дані дискретні. SMOTE та ADASYN зменшили PR-AUC бустингу з 0,179 до 0,12. SMOTENC зробив синтетичних пацієнтів класу 1 схожими на типових пацієнтів класу 0. Ваги класів та ROS лише зсувають поріг рішення і втричі погіршують калібрування ймовірностей. Той самий ефект без втрат дає вибір порога на крос-валідації.

\item Створені ознаки (сумарні візити, інтенсивність лікування, групи діагнозів, відхилення від середнього за спеціальністю) займають верхні позиції за взаємною інформацією, але покращують модель не більш ніж на одне стандартне відхилення. Вони також створюють мультиколінеарність, якої у вихідних даних не було.

\item Підсумковий конвеєр досяг на тестовій вибірці PR-AUC 0,187 проти 0,090 у випадкового класифікатора. Група ризику, яку він виділяє, повертається до лікарні в 2,3 раза частіше за середнє.
\end{enumerate}

Головний методологічний висновок роботи: методи обробки й аугментації слід оцінювати за якістю моделі на реальних даних, а не за їхніми власними критеріями. Наближення до нормального розподілу знищило інформацію про кількість візитів, вирівнювання класів зруйнувало калібрування ймовірностей, а синтетичні дані виявилися легко відмінними від реальних. Кожну з цих проблем виявила лише перевірка на моделі, і саме уникнення таких кроків дало найбільший внесок у якість.

\section*{Список використаних джерел}

\begin{enumerate}[label={[\arabic*]},ref={\arabic*},leftmargin=*,labelindent=0pt]
\item \label{ref:strack} Strack B., DeShazo J. P., Gennings C. et al. Impact of HbA1c Measurement on Hospital Readmission Rates: Analysis of 70,000 Clinical Database Patient Records // BioMed Research International. 2014. Vol. 2014. Article ID 781670.
\item \label{ref:saito} Saito T., Rehmsmeier M. The Precision-Recall Plot Is More Informative than the ROC Plot When Evaluating Binary Classifiers on Imbalanced Datasets // PLoS ONE. 2015. Vol. 10, No. 3. e0118432.
\item \label{ref:bergsma} Bergsma W. A bias-correction for Cramér's V and Tschuprow's T // Journal of the Korean Statistical Society. 2013. Vol. 42, No. 3. P. 323--328.
\item \label{ref:iforest} Liu F. T., Ting K. M., Zhou Z.-H. Isolation Forest // Proc. 8th IEEE International Conference on Data Mining. 2008. P. 413--422.
\item \label{ref:lof} Breunig M. M., Kriegel H.-P., Ng R. T., Sander J. LOF: Identifying Density-Based Local Outliers // Proc. ACM SIGMOD International Conference on Management of Data. 2000. P. 93--104.
\item \label{ref:mice} van Buuren S., Groothuis-Oudshoorn K. mice: Multivariate Imputation by Chained Equations in R // Journal of Statistical Software. 2011. Vol. 45, No. 3. P. 1--67.
\item \label{ref:target} Micci-Barreca D. A Preprocessing Scheme for High-Cardinality Categorical Attributes in Classification and Prediction Problems // ACM SIGKDD Explorations Newsletter. 2001. Vol. 3, No. 1. P. 27--32.
\item \label{ref:yeojohnson} Yeo I.-K., Johnson R. A. A New Family of Power Transformations to Improve Normality or Symmetry // Biometrika. 2000. Vol. 87, No. 4. P. 954--959.
\item \label{ref:smote} Chawla N. V., Bowyer K. W., Hall L. O., Kegelmeyer W. P. SMOTE: Synthetic Minority Over-sampling Technique // Journal of Artificial Intelligence Research. 2002. Vol. 16. P. 321--357.
\item \label{ref:adasyn} He H., Bai Y., Garcia E. A., Li S. ADASYN: Adaptive Synthetic Sampling Approach for Imbalanced Learning // Proc. IEEE International Joint Conference on Neural Networks. 2008. P. 1322--1328.
\item \label{ref:batista} Batista G. E., Prati R. C., Monard M. C. A Study of the Behavior of Several Methods for Balancing Machine Learning Training Data // ACM SIGKDD Explorations Newsletter. 2004. Vol. 6, No. 1. P. 20--29.
\item \label{ref:elor} Elor Y., Averbuch-Elor H. To SMOTE, or not to SMOTE? // arXiv preprint arXiv:2201.08528. 2022.
\end{enumerate}

\end{document}
